\section{Knowledge Types in Regulatory Compliance}

Knowledge can be broadly classified into three categories: explicit, implicit, and tacit knowledge.

\subsection{Explicit Knowledge}

Explicit knowledge refers to information that can be formally documented, stored, and transferred. In the compliance domain, this includes regulations, regulatory guidance, policies, procedures, and control frameworks. Such knowledge represents the codified requirements that AI systems can retrieve, process, and analyse.

\subsection{Implicit Knowledge}

Implicit knowledge refers to knowledge that can be inferred from relationships and patterns within explicit information. In regulatory compliance, obligations often span multiple documents, where relationships between regulations, amendments, supervisory guidance, and internal controls reveal additional meaning that is not directly stated within any single source.

\subsection{Tacit Knowledge}

Tacit knowledge represents experiential knowledge that is difficult to formalise or encode. Within regulatory compliance, it reflects the expertise required to understand how regulatory artefacts relate to one another, including knowing which related documents should be considered, when they should be reviewed, and how changes in one document may affect the interpretation of others.

The regulatory landscape can therefore be represented as a dynamic knowledge graph, where regulations, amendments, technical standards, and guidance documents form interconnected nodes linked through relationships such as ``amends'', ``references'', and ``clarifies''. While explicit knowledge represents the information contained within these nodes, tacit knowledge represents the ability to navigate the relationships between them. For AI systems, capturing this associative knowledge is essential for moving beyond document retrieval towards expert-level regulatory reasoning.